\reportmodule{1}{一}{选题报告}\label{chap:selection-report}

\section{选题依据}

\subsection{研究来源}

本课题来源于工具调用型大语言模型智能体的实际开发工作。智能体完成一项任务时，通常需要多次请求模型、调用外部工具并更新运行状态。任务失败后，开发者需要查阅提示、调用参数、工具返回、状态变化和错误信息，判断问题发生的环节及其对后续步骤的影响。这些材料分散在调用平台、日志、配置和代码中，增加了查找和理解的难度。

\subsection{研究背景}

大语言模型（large language model，LLM）通过工具调用参与外部任务执行，形成大语言模型智能体（large language model agent）。Yao 等\cite{yao2023react}提出的推理与行动协同方法（ReAct）将语言推理与环境操作相结合，使模型能够根据任务进展继续采取行动。Cemri 等\cite{cemri2025mast}对多智能体失败的分析以及 Chen 等\cite{ref20}对智能体框架缺陷的研究表明，故障可能涉及模型输出、工具交互、状态传递和流程控制等多个环节。开发者需要把最终症状与运行事件联系起来，比较原因解释并决定后续检查。

近期研究开始提供这一问题的直接证据。Epperson 等\cite{ref02}通过开发者访谈和用户研究发现，长对话检查、问题定位及运行调整存在困难，提供编辑和重执行功能后，使用者仍需判断修改位置和方式。van der Maden 等\cite{ref01}对十九名大语言模型产品实践者的访谈揭示，发现评价问题与采取具体改进行动之间存在缺口。Kumar 等\cite{ref59}对开发者与编程智能体协作的观察，也体现了持续评价、增量交互和人工干预在真实任务中的作用。

本课题研究开发者形成诊断判断的过程，重点关注运行材料的组织和使用。信息表征指运行事件及其属性、关系和来源在界面中的呈现方式；交互设计指开发者查找、关联、比较和核验这些信息的操作。研究将从实际诊断事件出发，确定具有共同性的困难，再形成和检验具体设计。

\subsection{研究目的及意义}

本研究拟建立工具调用型智能体故障诊断的信息条件与开发者活动之间的经验联系，围绕一项经研究确认的核心困难，探索能够改善诊断判断的信息表征与交互方式，并提炼具有适用条件的设计原则。

具体目标包括：

\begin{enumerate}
\item 通过真实事件访谈和诊断任务观察，描述开发者查找材料、形成原因解释和检验判断的过程，识别反复出现的困难及其条件。
\item 根据经验研究确定设计问题，比较备选方案，形成能够表达和检验核心设计主张的原型。
\item 通过对照任务和新案例观察，评价设计对诊断判断、证据使用和操作成本的影响，明确原则的作用与适用范围。
\end{enumerate}

\subsubsection{理论意义}

本研究面向交互设计领域，拟形成经验描述和设计原则两类学术成果。经验描述用于说明开发者在什么任务和信息条件下难以定位相关事件、建立事件关系、区分原因解释或决定后续检查，补充程序调试研究在自然语言记录、概率性输出和外部工具混合情境中的经验材料。

设计原则概括针对共同诊断困难的设计操作、预期作用和适用条件。Höök 和 Löwgren\cite{hook2012strong}讨论的中层知识（intermediate-level knowledge）为这一贡献定位提供了参考。原则将具体实例与可迁移的设计认识联系起来，使其他设计者能够判断其适用性并在不同界面中实施。原则由形成性研究、方案比较和评价结果共同形成，按照“问题条件—设计操作—预期作用—经验依据—适用边界”的结构表述。研究原型用于把设计主张转化为可操作方案，并在诊断任务中接受检验。

\subsubsection{实践意义}

本研究将形成可复现的诊断案例、研究原型和评价材料，为智能体调试工具的设计提供依据。研究结果可帮助工具设计者判断哪些信息需要被记录，哪些关系需要被呈现，以及哪些操作有助于使用者检查自己的判断。案例与评分材料还可为开发团队评价调试界面和改进运行记录提供参考。

\section{国内外研究综述}

\subsection{工具调用型智能体故障诊断}

\subsubsection{智能体执行机制}

大语言模型（large language model，LLM）能够根据自然语言目标生成操作方案，并调用外部工具执行任务。Yao 等\cite{yao2023react}提出的推理与行动协同方法（ReAct）将语言推理和环境操作组织为交替过程，为大语言模型智能体（large language model agent）提供了一种基本实现方式。Cemri 等\cite{cemri2025mast}对多智能体失败的分析以及 Chen 等\cite{ref20}对智能体框架缺陷的研究均表明，失败分析涉及模型输出、工具交互、状态传递和流程控制。

\subsubsection{故障诊断任务}

Epperson 等\cite{ref02}的开发者研究记录了长轨迹检查、干预位置选择和运行调整中的困难；Wu 等\cite{ref23}则将任务、工具、数据和人工干预组织为可追溯的图结构，以支持失败分析。基于这些研究，本课题将故障诊断界定为开发者从失败症状出发，查找相关运行事件、形成原因判断并通过后续检查加以核验的过程。信息表征指运行事件及其属性、关系和来源在界面中的组织与呈现；交互设计指开发者进入、查找、比较和核验这些材料的操作方式。

\subsubsection{以人为中心的人工智能}

国内人机交互研究为这一问题提供了学科背景。范俊君等\cite{ref42}讨论了智能时代交互模型、界面与设计原则的发展，强调人的能力与机器能力在交互中的结合；许为、葛列众\cite{ref63}将人机合作、决策控制和以人为中心的人工智能纳入工程心理学研究框架。这些工作提出了研究方向。具体到智能体排错，仍需通过开发者的实际诊断行为，说明界面应支持什么活动，以及设计效果如何被观察和评价。

\subsection{程序调试的信息行为}

\subsubsection{信息觅食理论}

程序调试研究较早就把搜索位置与理解原因作为两个相互联系的活动。Katz 和 Anderson\cite{ref15}通过四项程序调试实验分析学习者的缺陷定位策略，发现搜索策略会影响检查路径，而定位到某一位置之后，仍然需要判断局部行为是否正确。该研究的对象是程序学习者，其价值在于把“查到哪里”和“如何判断”分开考察。对于智能体运行记录，快速到达一次异常调用只是诊断过程的一部分，开发者还需解释该调用与任务失败之间的联系。

Ko 等\cite{ref14}观察十名开发者完成软件维护任务，将其活动概括为查找、关联和收集相关信息。开发者需要在代码位置、依赖关系和临时收集的材料之间往返，信息线索不足及导航成本会妨碍理解。该研究将材料访问和信息保持确定为程序理解中的具体问题。本研究据此考察开发者怎样在提示、参数、工具返回和状态记录之间查找并关联材料。

Lawrance 等\cite{ref08}以信息觅食（information foraging）视角分析专业程序员的调试行为，研究信息线索、预期价值和获取成本如何影响下一步搜索。其解释重点在于开发者为何沿某条路径继续查找，以及什么时候离开当前材料。界面中的标签、摘要和跳转关系由此构成可观察的搜索线索，其作用需要结合实际访问路径和诊断结果评价。

这些研究共同建立了一组可观察的行为：进入材料、选择线索、建立关联、保留发现和判断原因。由于单独的访问次数和停留时间不能说明诊断是否有效，本研究将操作记录、参与者当时的问题和最终判断结合分析。

\subsubsection{程序理解}

Starr 和 Storey\cite{ref16}通过对二十七名专业开发者的访谈，运用建构主义扎根理论（constructivist grounded theory）研究排错中的认知经历。他们将理解异常行为的过程与困惑、心理模型构建以及认知疲劳联系起来。基于该研究，本研究记录参与者的编程经验、智能体开发经验和案例熟悉程度，并为正式任务提供统一的系统背景，以控制前置理解差异。

Romero 等\cite{ref34}研究多表征软件环境中的调试策略，发现表征使用与编程经验、表征知识及具体任务有关。本研究据此将表征理解、跨视图操作和参与者经验纳入原型走查与正式评价。

\subsubsection{意义建构理论}

意义建构（sensemaking）关注人如何把分散信息组织为能够理解情境和指导行动的解释。Pirolli 和 Card\cite{ref17}基于认知任务分析与口语报告，描述情报分析中的信息搜集和解释形成活动，并讨论技术介入的位置。其研究提供了从材料选择到解释形成的过程视角，但具体活动结构需要结合智能体诊断资料确定。

Klein 等\cite{ref11}提出的宏观认知模型强调资料与解释框架之间的相互作用：已有理解影响人注意哪些信息，新信息又可能促使人保留、补充或调整理解。该模型可用于分析原因假设如何随诊断材料发生变化，分析对象包括假设提出、材料选择、解释修订和后续检查。

Dritsa 和 Houben\cite{ref57}通过设计研究者使用健康相关数据可视化的经验研究，提出关注不确定性的意义建构模型，讨论数据、表征和使用者背景如何影响洞见形成。该研究展示了依据实践资料修订过程模型的方法，也说明理论框架需要通过具体活动资料检验。

本研究将信息觅食理论用于分析材料入口、搜索线索和访问成本，将意义建构理论用于分析原因假设的形成与修订。两种理论在形成性研究中的适用范围依据实际资料确定。

\subsubsection{自动故障定位}

自动故障定位（automated fault localization）根据程序运行信息生成可疑位置或排序。Parnin 和 Orso\cite{ref35}通过用户研究考察自动故障定位技术对程序员的实际帮助，指出提供可疑代码排序之后，开发者仍需要上下文才能理解和修复问题。该研究把算法提供的线索与人完成的诊断工作联系起来，也表明评价自动定位工具时，需要观察使用者能否把提示转化为可执行的检查。

自动生成的原因解释属于诊断输入，其实际作用取决于开发者能否理解和核验建议依据。本研究将“取得线索—解释异常—检验判断”作为初始分析框架，并通过形成性研究检验这些活动在智能体任务中的具体表现。

\subsection{大语言模型开发中的人机协作}

\subsubsection{生成式人工智能评价}

Barke 等\cite{ref64}观察二十名程序员使用代码生成工具，并通过扎根分析描述加速既定实现和探索可能方案两种交互方式。当开发者已有明确意图时，生成结果可以加快实现；当其探索陌生任务时，需要通过阅读、试用和修改来建立理解。该研究提示，人使用生成系统的目标和知识状态会改变其信息需要。

Zamfirescu-Pereira 等\cite{ref65}研究非人工智能专家设计提示的过程，发现参与者会基于有限输出进行机会式修改，并将日常人际交流经验带入对模型的判断。该研究的参与者主要是非人工智能专家，其结论尚不能直接说明专业开发者的诊断行为。

van der Maden 等\cite{ref01}访谈十九名大语言模型产品实践者，分析评价工作如何连接产品改进。研究发现，识别质量问题、解释问题来源和决定修改措施之间存在可行动性缺口。本课题据此关注失败被识别之后的原因判断和检查决策。

以上研究分别涉及代码生成、提示设计和产品评价，共同指向一种持续存在的工作：使用者需要判断生成结果的含义及其对下一步行动的影响。智能体增加了多步骤执行和外部工具反馈，因此需要进一步研究，开发者究竟依据哪些事件建立解释，如何排除相似原因，以及何时决定补充观测。

\subsubsection{编程智能体协作}

Kumar 等\cite{ref59}研究十九名开发者与编程智能体协作处理三十三个真实仓库问题的过程。参与者处理自己曾贡献过的仓库任务，研究观察到增量交互、主动干预和持续评价在协作中的作用。其材料保留了任务意图、仓库约束和既有知识，能够揭示行为模式及关联，具体交互方式的因果效果仍需控制条件下的比较。本研究参照其取样思路，先从真实事件中识别诊断困难，再使用可复现系统构建正式任务。

\subsubsection{自动化透明度}

自动化透明度（automation transparency）关注系统过程和状态信息如何向使用者呈现。Pareek 等\cite{ref58}结合文献分析与实验室研究，让十二名参与者使用五种多智能体界面完成信息查找和逻辑推理任务，并分别收集透明性、帮助程度和可靠性判断。该研究说明主观评价包含多个维度，本研究据此将主观感受与诊断表现、证据使用和判断信心分别测量。

\subsection{智能体调试的交互设计}

\subsubsection{执行轨迹可视化}

Cornelissen 等\cite{ref33}通过受控实验研究执行轨迹可视化对程序理解的影响，在其任务设置下观察到正确性和完成时间方面的改善。该实验把可视化效果与具体程序理解任务相联系，为诊断界面采用任务表现和完成时间指标提供了依据。

Ko 和 Myers\cite{ref07}开发面向原因问答的调试界面（Whyline），让开发者从可见程序输出提出“为什么发生”与“为什么没有发生”的问题，再进入相关执行事件。该研究通过用户任务检验了问题入口与执行事件连接的作用。本研究将从失败症状进入相关运行记录的操作作为候选设计方向，其必要性由形成性研究确定。

上述研究处理的执行记录具有明确的程序语义。工具调用型智能体的运行材料同时包含调用结构、数据传递和自动分析结果。本研究将在案例材料中标记关系来源，并在原型评价中检查参与者对关系语义的理解。

\subsubsection{结构化复盘}

Khanna 等\cite{ref04}通过人工智能复盘工具（AAR/AI）的实证研究，考察结构化反思如何帮助使用者发现人工智能系统的缺陷。研究将相同解释材料置于不同组织条件下，以“发生了什么、为什么、需要改变什么”等问题支持复盘，并测量故障回忆和描述。结果为信息组织方式影响缺陷理解提供了证据。其研究对象和任务具有特定背景，应用于大语言模型智能体时，需要检验结构化问题是否符合开发者的诊断活动。

Arawjo 等\cite{arawjo2024chainforge}提出提示与输出比较工具（ChainForge），通过可视化流程支持提示、模型和输入的组合试验。其评价包括二十一人的实验室研究，以及六次涉及八名使用者的访谈，揭示了工具在输出探索、假设检验和评价迭代中的使用方式。该工作说明，将多个输出放在可比较的结构中，可以支持对模型行为的考察。

上述研究分别检验了结构化复盘和输出比较。本研究不预设比较对象，将依据形成性研究中出现的诊断任务确定比较内容，并通过任务表现、操作记录和使用成本评价其作用。

\subsubsection{分层可视化}

Lu 等\cite{ref22}提出智能体行为可视分析系统（AgentLens），通过层级结构和多粒度时间信息支持对自主智能体行为的探索。论文报告了使用案例和十四名参与者的用户研究，体现了从概览进入局部行为的设计价值。该系统主要处理行为分析问题，研究材料能够说明使用者如何理解复杂运行过程；其结果还需要结合具体故障答案，才能支持关于诊断正确性的结论。

Gao 等\cite{ref24}提出科学智能体执行轨迹可视化工具（Graph of Trace），将运行过程组织为可交互的事件图，并请相关领域专家评价系统。专家反馈提供了可理解性和使用价值方面的证据。科学任务中，领域知识会影响对中间结果的解释，工具对专业使用者的帮助需要结合其任务背景理解。

这些工作展示了概览、分层和关系连接的实现方式，但其评价主要集中于可理解性和使用价值。时间顺序、调用包含、数据依赖和原因假设具有不同语义，本研究将关系识别和局部材料访问纳入原型评价。

\subsubsection{交互式调试}

Dibia 等\cite{ref25}提出多智能体开发工具（AutoGen Studio），整合组件配置、执行过程查看和调试功能。该系统论文展示了开发流程与工具能力，为案例组织和基础界面提供了参考。其部署经验有助于识别常见操作需求，具体设计对人的诊断效果仍需相应用户任务数据。

Epperson 等\cite{ref02}通过五名开发者的形成性访谈提出交互调试工具（AGDebugger），随后分别开展诊断和运行引导研究。诊断研究涉及六名参与者，运行引导研究涉及八名参与者。参与者对编辑、重置和局部重执行等功能表现出偏好，但研究也记录了判断修改位置和修改方式的困难。本研究将在形成性研究中考察重执行是否属于核心诊断活动；若纳入原型，则记录干预内容、受控条件和结果变化。

Rorseth 等\cite{ref28}提出面向数据应用的智能体调试器（LADYBUG），展示逐步追踪、局部干预及受影响环节重新执行的实现。Hutter 和 Pradel\cite{ref18}提出软件开发智能体调试器（AgentStepper），提供暂停、检查和编辑等功能，并报告小规模学生参与者研究。前者属于系统演示，后者提供了小规模任务评价，尚未形成专业开发者在多类故障上的比较证据。

\subsubsection{诊断工具评价}

\begingroup
\small
\setstretch{1.15}
\renewcommand{\arraystretch}{1.1}
\setlength{\LTleft}{0pt}
\setlength{\LTright}{0pt}
\begin{longtable}{@{}>{\raggedright\arraybackslash}p{0.23\dimexpr\textwidth-6\tabcolsep\relax}>{\raggedright\arraybackslash}p{0.26\dimexpr\textwidth-6\tabcolsep\relax}>{\raggedright\arraybackslash}p{0.24\dimexpr\textwidth-6\tabcolsep\relax}>{\raggedright\arraybackslash}p{0.27\dimexpr\textwidth-6\tabcolsep\relax}@{}}
\caption{诊断工具研究比较}\label{tab:literature-review-1}\\
\toprule
\textbf{研究} & \textbf{主要研究对象与方法} & \textbf{已提供的证据} & \textbf{对本课题仍需检验的内容} \\
\midrule
\endfirsthead
\multicolumn{4}{c}{表\thetable\quad 诊断工具研究比较（续）}\\
\toprule
\textbf{研究} & \textbf{主要研究对象与方法} & \textbf{已提供的证据} & \textbf{对本课题仍需检验的内容} \\
\midrule
\endhead
\bottomrule
\endfoot
Whyline\cite{ref07} & 程序输出原因查询与用户任务评价 & 问题入口与执行事件连接的价值 & 自然语言记录及不完整观测中的关系理解 \\
AAR/AI\cite{ref04} & 结构化复盘的对照研究 & 相同解释材料的组织方式影响缺陷描述 & 结构化复盘在专业智能体排错中的作用 \\
ChainForge\cite{arawjo2024chainforge} & 实验室研究与实际使用访谈 & 输出比较、假设检验和评价迭代的使用方式 & 跨步骤故障解释与下一步检查 \\
AgentLens\cite{ref22} & 使用案例与用户研究 & 层级轨迹支持复杂行为探索 & 诊断正确性、原因判断及证据使用 \\
AGDebugger\cite{ref02} & 开发者访谈、诊断与运行引导研究 & 编辑和重执行的需求及使用困难 & 哪种信息支持使用者选择干预位置 \\
AgentStepper\cite{ref18} & 小规模任务研究 & 单步检查与编辑的初步可用性证据 & 经验差异及多类案例上的效果 \\
Graph of Trace\cite{ref24} & 科学智能体案例与专家评价 & 事件图的可理解性与使用价值 & 主观理解和客观诊断表现之间的关系 \\
\end{longtable}
\endgroup

现有工具研究采用了问题入口、结构化复盘、输出比较、层级轨迹和局部重执行等设计方式。其评价对象、任务和样本差异较大，尚不足以比较这些方式对专业开发者故障诊断的作用。本研究将依据形成性研究选择核心设计主张，并通过可比较任务检验其效果。

\subsection{智能体运行追踪}

\subsubsection{分布式追踪}

分布式追踪（distributed tracing）研究为跨组件运行信息的关联提供了技术基础。Fonseca 等\cite{ref30}提出跨层任务追踪框架（X-Trace），通过传播任务信息重建请求经过的组件与操作。Sigelman 等\cite{ref32}总结大规模追踪基础设施（Dapper）的部署经验，讨论追踪开销、采样及跨服务关联。Mace 等\cite{ref31}提出动态因果监测系统（Pivot Tracing），通过动态插桩和查询支持跨组件分析。这些工作说明，统一标识与关系记录能够将分散事件组织为可查询的执行过程。

分布式追踪研究表明，事件关联取决于统一标识、上下文传播和数据采集范围。本研究将先依据诊断任务确定所需材料，再核查案例中的记录能力；无法由原始记录支持的关系将标明其推断来源。

Zheng 等\cite{ref19}提出系统级智能体观测工具（AgentSight），关联模型交互与底层系统活动，并在异常行为和性能分析场景中展示其作用。该研究补充了应用层对话记录之外的信息来源。采集层级的扩展能够帮助解释某些运行问题，同时也会增加材料种类；开发者怎样选择和组合这些信息，仍然是交互研究需要回答的问题。

\subsubsection{自动故障归因}

Cemri 等\cite{cemri2025mast}分析多智能体系统的失败运行并建立分类，涉及系统设计、智能体间协调以及任务验证等问题。该分类有助于理解智能体失败并不限于单次模型输出，也可作为案例采样的参考。多智能体协作中的协调和责任分配包含额外变量，因此借鉴其故障类型时，应逐项判断是否适用于单个工具调用型智能体。

Zhang 等\cite{ref03}围绕“哪个智能体在何时导致失败”构建自动失败归因任务，并比较不同归因方法。研究表明，在已有轨迹中定位责任主体和责任步骤仍具有挑战。这类评测回答机器方法能否接近标注答案；对于人使用诊断工具的研究，还需检查参与者是否理解建议依据、是否能够识别错误建议，以及是否会把局部责任步骤视为完整原因。

Chen 等\cite{ref20}研究智能体运行框架缺陷的诊断和修复，将失败轨迹与数据流、控制流及实现位置联系起来。研究提供了从运行表现追查框架问题的技术路径，其评价重点是自动诊断和修复。轨迹与实现的关联可以成为界面候选功能，但其必要性需要根据开发者的具体任务判断：有些问题可以在工具输入输出层面解释，有些则需要继续进入配置或实现细节。

上述研究为故障案例构建提供分类和实现依据，也表明自动建议需要保留来源、模型、输入和生成版本。本研究将自动归因结果与原始运行事实分别记录，用于分析使用者怎样接受、质疑或放弃建议。

\subsubsection{运行关系可视化}

Wu 等\cite{ref23}提出轨迹转图分析平台（AgentGraph），将智能体、任务、工具、输入输出和人工干预组织为具有类型的节点与关系，并将图元素链接到原始轨迹片段。系统还展示了失败分析和扰动测试。该工作说明，图结构可以同时承载运行对象、关联信息和原始出处。

AgentGraph 的图元素同时包含调用关系、数据关系和分析结果。由于这些关系的证据来源不同，本研究将来源标识、关系说明和原始材料访问作为图结构方案的评价内容。

\subsubsection{来源追溯}

来源追溯（provenance）研究关注数据与分析结果从何而来。Ragan 等\cite{ref36}提出可视化和数据分析中的来源类型与用途框架，区分数据、操作、交互、洞见和理由等记录。该区分提示，智能体执行历史与开发者的分析历史承担不同作用：前者描述系统发生的事件，后者描述使用者如何形成判断。研究过程需要同时保留两者，才能将界面操作与诊断解释联系起来。

Kadivar 等\cite{ref38}提出可编辑和可回放分析历史的可视分析系统（CzSaw），支持使用者回顾、调整和复用分析过程。Xu 等\cite{ref37}进一步讨论分析来源如何服务于理解、审计和协作，并提出将操作记录与推理活动连接的研究议程。这些工作为记录已检查材料、保留暂定假设和比较分析分支提供了设计资源。

来源追溯研究区分系统执行记录和使用者分析记录。本研究据此在操作记录之外采集参与者的阶段性判断和理由，用于分析运行材料如何进入诊断结论。

\subsection{可解释人工智能}

\subsubsection{解释需求}

可解释人工智能（explainable artificial intelligence，XAI）研究将解释放在人使用系统的情境中考察。Liao 等\cite{ref49}访谈二十名相关设计实践者，并借助问题提示讨论解释需求，形成面向用户问题的设计资源。研究显示，使用者需要的解释与其角色、任务及所处环节有关。本研究据此区分故障理解、修改位置选择和补充观测等解释用途。

Ehsan 等\cite{ref51}以情境化设计材料访谈二十九名人工智能使用者与实践者，研究社会透明性（social transparency），将其他人的使用情境、行动和结果纳入解释。该研究表明，解释需求包含技术输出之外的工作背景。本研究将在真实事件访谈中记录配置变更、既往检查和协作信息，再依据跨案例分析决定是否纳入设计。

基于上述研究，形成性访谈将围绕具体诊断事件询问参与者试图确认的问题、使用的材料和采取行动的依据，避免仅收集脱离任务情境的功能偏好。

\subsubsection{解释使用偏差}

Kaur 等\cite{ref50}结合十一名数据科学家的情境研究与一百九十七人的调查，考察解释工具在机器学习分析中的实际使用。研究观察到使用者对解释结果的误解和过度依赖。该发现表明，解释功能的使用和理解需要分别评价。

Buçinca 等\cite{ref55}在一百九十九名参与者的人工智能辅助决策实验中比较多种交互条件，发现促使使用者独立思考的认知强制措施（cognitive forcing functions）可以减少对错误建议的过度依赖，但也伴随接受度和使用成本方面的变化。该研究据此分别报告界面偏好、操作时间和判断正确性。

Sieker 等\cite{ref56}在视觉问答任务中操纵系统可见的图像信息，考察解释是否帮助使用者理解系统能力与局限。研究中，提供解释提高了使用者对系统能力的评价，却未帮助其更准确地识别系统限制。使用者是否感觉“理解了”，需要与其能否识别系统限制及错误结果一起评价。

三项研究共同表明，主观理解、判断正确性、使用时间和对系统建议的依赖需要分别测量。本研究据此将诊断表现、证据使用、判断信心和使用体验设为不同评价指标。

\subsubsection{不确定性可视化}

Hullman\cite{ref40}调查九十名可视化作者并访谈十三名设计者，研究实际传播中为何常省略不确定性。其结果揭示了信息简洁性、沟通目标、受众预期和表达成本之间的张力。研究对象是可视化作者，本研究将通过开发者任务另行考察不确定性信息对后续检查的作用。

Guo 等\cite{ref62}通过用户实验比较图边不确定性与其他边属性的配对视觉编码，发现编码效果受到视觉变量组合和任务类型影响。本研究因此在具体诊断任务中评价不确定性编码，检查使用者能否区分数据缺失、关系不确定和推断强度。

Sacha 等\cite{ref39}从可视分析过程讨论不确定性、意识与信任的关系，为分析不确定性在处理和理解环节中的传播提供概念框架。Ehsan 等\cite{ref54}通过四十三名参与者的情境化协同设计，研究如何在解释中显露系统接缝、限制和工作条件。前者提供分析视角，后者提供设计探索的经验材料，共同支持把系统限制作为可供判断的信息。

结合来源追溯和不确定性研究，本研究在案例构建中区分三类信息状态：运行事实已经记录但存在多种解释、关键事实缺失、自动建议尚未核验。该分类用于组织初始案例，并根据形成性研究资料修订。

\subsection{通过设计开展研究}

\subsubsection{设计原则的形成}

Amershi 等\cite{ref48}通过多轮形成、评议和验证提出十八条人机交互指南，评价涉及四十九名设计实践者和二十种产品。该研究展示了通用指导如何通过多种产品和使用情境检验其清晰性与适用性。原则的价值既体现在表述，也体现在其能否帮助设计者识别问题和作出设计决策。

Weisz 等\cite{ref66}结合文献、设计实践者反馈、现有生成式人工智能应用以及两个应用的设计过程，迭代形成六项设计原则及相应策略。研究把原则的抽象表述与具体实施方式联系起来，并通过多种来源修订其内容。上述两项研究表明，设计原则的形成需要经验问题、设计操作、任务结果和应用条件相互对应。

\subsubsection{中层知识}

Zimmerman 等\cite{zimmerman2007rtd}讨论通过设计开展研究（research through design）的知识生产方式，将设计实践与研究贡献联系起来。Höök 和 Löwgren\cite{hook2012strong}以强概念（strong concepts）讨论交互设计中的中层知识（intermediate-level knowledge），说明知识可以位于单个实例与宏观理论之间。两项工作支持设计研究提炼可迁移认识，同时要求研究者明确其知识形态及依据。

本课题的设计学贡献定位为具有条件的设计原则，以及支撑这些原则的跨案例经验描述。研究原型承担设计主张的物化和检验功能，形成性研究、方案比较和实证评价共同构成设计知识的生成过程。原则应说明面对哪类诊断困难，设计者可以怎样组织信息或提供操作，预期影响什么活动，现有证据来自哪些参与者和任务，以及何时效果减弱或成本增加。原则可以借助一个原型检验，但其表述需要让其他设计者在另一种界面中实施。

\subsubsection{研究严谨性}

Sedlmair 等\cite{sedlmair2012design}总结可视化设计研究的方法，将问题理解、抽象、设计实现和验证联系起来，并讨论各阶段的研究风险。Meyer 和 Dykes\cite{meyer2020rigor}进一步从可视化设计研究的经验出发讨论严谨性，强调研究过程的反思、论证和材料呈现。本研究依据这些方法记录设计选择、修改理由和对应证据。

候选原则需要追溯到跨案例困难、可比较方案和评价结果。实验条件同时改变多个功能时，研究结论限定为组合设计的整体效果；单项功能的作用需要相应的条件比较或过程证据。

\subsubsection{实证评价方法}

Malterud 等\cite{malterud2016information}提出信息力（information power），将定性样本的充分性与研究目标、样本特征、理论使用、访谈质量和分析方式联系起来。该观点支持以研究问题和资料质量决定招募，而不是仅以访谈人数判断充分性。对于本课题，关键在于是否获得不同开发背景下可分析的诊断事件，以及这些事件能否支持跨案例比较。

Gale 等\cite{gale2013framework}介绍框架法（framework method），通过编码和矩阵组织案例与主题，适用于系统比较多名参与者的经验。本研究将建立案例—主题矩阵，并保留不符合主要模式的案例，用于区分跨案例困难、特定任务困难和个体差异。

对照实验中的样本量需要另行论证。Lakens\cite{ref67}讨论最小有意义效应、估计精度和研究资源等样本量依据；Green 和 MacLeod\cite{ref68}提供广义线性混合模型的模拟功效分析方法。参与者反复处理多个案例时，人员和案例都可能造成差异，样本论证应纳入这种结构。预实验适合用于检查任务难度、测量流程和参数范围，正式比较应在明确主要指标和分析计划之后开展。

\subsection{研究综述小结}



现有文献形成了四方面认识。第一，调试包含材料搜索、关系理解、假设形成与检查，自动提供可疑位置仍需要人的解释工作。第二，生成式系统的实践者持续面临结果评价和改进行动之间的衔接问题，智能体增加了跨步骤执行、工具反馈和干预选择。第三，轨迹可视化、问题式入口、输出比较和局部重执行提供了多种设计资源，效果与任务、使用者经验及信息条件有关。第四，清晰解释和透明性感受需要与客观判断分别评价，不确定性和系统限制会影响使用者的检查行为。

国内研究在智能时代人机协作与工程心理学方面形成了重要议题，国际相关文献积累了程序调试、生成式工具使用和智能体可视分析的具体经验。现阶段，工具调用型智能体的专业开发者诊断行为，尤其是运行材料怎样影响原因判断和下一步检查，仍需要更直接的经验研究。



首先，需要补充从真实诊断事件出发的过程证据。现有工具评价揭示了使用需求和部分任务表现，但关于开发者如何在工具参数、返回内容、状态变化和配置之间寻找原因，仍缺少足以支撑设计选择的细致描述。研究应确定困难出现的条件，并区分材料未被记录、材料难以访问以及材料虽已取得但难以解释。

其次，需要补充针对核心设计选择的比较证据。多种功能组合可以改善工具体验，但研究还需说明设计改变了哪项诊断活动。可在控制材料与基础功能的情况下比较信息组织方式，也可明确研究“新增信息及其呈现”的整体效果。评价对象的界定决定结论能否解释具体设计的作用。

再次，需要补充对判断依据和材料边界的评价。只按最终原因名称评分无法区分有证据支持的判断、无依据猜测和材料不足时的合理保留。本研究将实际故障原因和可见材料允许作出的判断分别用于案例记录和评分。

最后，需要补充设计知识在新案例中的适用证据。新案例可以改变工具组合、任务内容、故障传播方式或记录完整性，研究据此检查某项原则是否依赖具体案例或界面实现。跨案例验证能够逐步明确知识适用的情境，支持后续设计者判断能否借鉴。



本课题聚焦具有工具调用和状态更新过程的单个主要智能体，研究开发者在任务失败后如何利用运行材料形成并检验诊断判断。研究首先通过访谈和任务观察建立跨案例描述，再选择一项具有共同性且能够通过设计改善的困难，形成可比较的原型方案，随后开展对照评价与新案例检验。

核心研究问题为：在工具调用型智能体故障诊断中，哪些信息条件使开发者难以建立和检验原因判断，以及怎样的信息表征与交互支持能够改善这些活动？研究拟形成诊断困难与信息条件的经验描述，以及由实证结果支持的设计原则。原型用于检验设计主张，对照实验评价其作用和成本，新案例研究界定适用范围。

\section{研究问题与方法}

\subsection{研究问题}

本课题的核心问题是：\textbf{在工具调用型智能体故障诊断中，哪些信息条件使开发者难以形成和检验原因判断，怎样的信息表征与交互支持能够改善这些活动？}

围绕核心问题提出三个研究问题：

研究问题一(RQ1)：在智能体发生任务失败后，开发者如何利用运行信息判断故障原因，诊断过程中存在哪些反复出现的困难？

研究问题二(RQ2)：针对上述诊断困难，如何组织和呈现相关运行信息，并通过交互支持开发者形成和验证故障原因假设？

研究问题三(RQ3)：所提出的交互设计如何影响开发者的故障诊断结果与诊断过程，这种影响在不同故障案例中有何表现？

\subsection{研究对象}

研究对象为具有实际开发或排错经验的工具调用型智能体开发者。目标系统包含一个主要智能体，能够通过多次模型请求和工具调用完成任务，并保留可供检查的运行记录。研究任务从失败结果出发，要求开发者解释异常、指出依据并提出下一步检查。

主要分析单位为一次具有明确任务目标、失败表现和诊断活动的事件。案例将覆盖跨步骤信息传递、工具交互及流程控制等运行条件。多智能体研究为相关工作和案例设计提供参考，正式研究范围集中在单个主要智能体，以保持责任关系和任务条件的可分析性。

研究关注对运行行为的诊断，包括工具参数、返回内容、状态更新和配置条件。模型内部表征、训练机制和参数级解释不构成本课题的研究对象。

\subsection{理论基础}

本研究以任务活动和运行材料为分析起点，将任务症状、可见材料、参与者问题、原因假设、检查行动和判断变化作为初始分析概念，并在开放编码过程中补充和修订。

信息觅食视角用于解释材料入口、搜索线索和访问成本。意义建构（sensemaking）视角用于分析分散材料如何被组织为原因解释，以及解释怎样随新材料改变。Pirolli 和 Card\cite{ref17}的认知任务分析提供了信息搜集与解释形成的过程参照；Klein 等\cite{ref11}关于资料与解释框架相互作用的讨论，可用于识别假设维持与修订的条件。

资料分析分别编码材料访问、搜索线索、原因假设、判断修订和检查行动。研究保留不能由初始理论框架解释的资料，并据此调整分类。

\subsection{研究方法}

\subsubsection{形成性研究}

\noindent\textbf{参与者招募。}

形成性研究拟招募12—16名参与者，通过公司和领英（LinkedIn）等外部专业渠道招募，覆盖不同组织、开发经验和工具使用背景。纳入条件为：近六个月内实际搭建、修改或维护过工具调用型智能体，且能够描述至少一次亲自参与的排错事件。招募将记录开发年限、智能体经验、系统类型、排错频率和来源渠道，并针对已有样本的缺项补充参与者。

Malterud 等\cite{malterud2016information}提出的信息力（information power）将定性样本充分性与目标、样本特征、访谈质量和分析方式联系起来。本研究据此评估资料是否足以回答研究问题，重点检查核心困难是否有跨参与者和跨任务的支持，反例是否得到解释，以及不同经验背景的材料是否充分。计划人数用于安排研究资源，实际招募根据资料质量和案例覆盖进行调整。

\noindent\textbf{资料收集。}

每次研究约60—90分钟，包括关键事件回溯和诊断任务观察。关键事件回溯围绕参与者亲历的故障展开，依次询问任务目标、最初症状、已取得的材料、曾怀疑的原因、采取的检查和最终判断。优先请参与者借助获得授权并脱敏的记录说明过程；无法提供记录时，以事件时间线和具体操作回忆补充，并标记资料来源。

任务观察使用初步构建的失败案例，让参与者查阅运行材料并完成诊断。研究记录屏幕、操作和同步口语报告（think-aloud），在自然停顿处询问当前判断及其依据。任务结束后回看关键片段，核对研究者对操作含义的理解。真实事件用于理解工作情境，统一任务便于比较不同参与者在相近条件下的活动。

资料包括访谈转录、事件概要、运行片段、操作时间线、参与者提出的假设及判断变化。研究将区分参与者回忆的经历、任务中直接观察的行为和研究者的解释。

\noindent\textbf{资料分析。}

研究采用框架法（framework method）。Gale 等\cite{gale2013framework}提出的案例—主题矩阵适合组织多名参与者的过程资料。分析先选取不同背景的案例开放编码，再结合初步分析框架形成编码表，随后在其余资料中应用和修订。

编码内容主要包括：当前诊断问题、所需与实际取得的信息、事件关系、原因解释、检查行动、判断变化及活动结果。跨案例比较识别同一困难是否出现于不同系统或任务，并结合口语报告、操作记录和阶段性判断区分相似行为所对应的活动。

研究拟邀请另一名研究人员复核约四分之一的资料，比较编码依据并记录分歧处理。分析保留反例和条件差异，形成核心困难、支持案例及适用条件的矩阵。结果进入设计阶段前，将邀请部分参与者核对事件解释是否符合其经历。

\subsubsection{案例构建}

\noindent\textbf{案例来源。}

研究以开源智能体框架自行构建可复现案例，并依据实践研究修订其任务和故障条件。初步考虑资料检索、结构化数据查询和多步骤工具处理等任务，最终选择取决于访谈结果及实现可行性。公司案例主要用于理解实际困难和核查任务真实性，正式实验材料以可公开或可授权复现的系统为主。

拟建立12—16个候选案例，经过专家核查和预实验后选取12个进入正式案例池。每个案例保存任务说明、系统版本、工具定义、正常运行材料、失败运行材料、配置及故障构造记录。保留完成诊断所需的背景和合理干扰事件，控制内容长度与领域知识要求。

案例条件依据形成性研究中的高频困难和差异条件确定。正式实验选择一项与核心设计主张直接相关的案例条件作为主要比较因素，其余条件用于平衡任务难度。

\noindent\textbf{证据评定。}

每个案例建立两份记录。故障记录说明实际改变了哪些系统条件、故障如何发生，以及哪些复现或干预支持这一解释。任务记录说明参与者可看到哪些材料、能够据此作出什么判断，以及仍有哪些无法确定的内容。

材料充分的案例要求参与者指出与证据相符的原因及其作用过程。材料不足的案例允许保留多个解释，但要求说明具体缺少什么，以及哪项检查能够区分这些解释。研究将以有针对性的缺失构造后一类任务，并核查剩余线索是否意外泄露答案。

案例拟由两名具有相关开发经验的人员独立检查。核查内容包括情境合理性、故障可复现性、可见材料充分性、替代解释和评分规则。存在多个有效诊断路径时，将其纳入评分手册。构造时选定的故障只是实际原因记录，评分还需依据参与者可见材料判断答案是否合理。

\subsubsection{原型设计}

本阶段采用通过设计开展研究的方法。原型用于外化设计假设、组织任务材料并产生可分析的使用证据，其迭代过程与取舍记录构成设计知识提炼的依据。

研究围绕形成性研究确认的核心困难提出至少两种备选方案。每种方案记录对应问题、设计操作、预期作用和实现条件，具体信息结构和交互方式在形成性研究后确定。

设计探索拟开展两轮任务式走查，共招募6—8名参与者，每轮约4—5人，部分参与者参与两轮以观察修改效果，并补充首次使用者检查学习成本。走查使用相同或难度相近的任务，收集理解错误、检查行为、未满足的信息需要和操作障碍。

原型分为通用功能和研究功能。通用功能包括任务说明、事件查看、必要搜索和材料展开；研究功能承载核心设计主张。每次迭代保存对应困难、设计意图、观察结果及修改理由，淘汰仅受偏好支持但未改善目标活动的方案。

最终原型优先使用保存的运行材料，以保证任务可复现。若形成性研究确认交互重执行对核心问题不可缺少，将为相关工具提供固定响应或受控执行，并记录每次修改和执行条件。正式评价前冻结原型版本。

\subsubsection{对照实验}

\noindent\textbf{实验条件。}

正式研究采用被试内设计（within-subject design），比较基础界面与研究界面。基础界面提供可用的时间序列、事件详情、搜索和展开功能；研究界面在此基础上实现选定的信息组织或交互支持。训练使用独立练习案例，确保参与者能够操作两种界面。

比较前建立条件清单，逐项核对原始材料、派生内容、搜索能力、默认展开状态和访问步骤。若研究重点为信息组织，两种界面提供相同语义内容，差异集中于其呈现和操作。若选定方案引入新的派生信息，则把“派生信息及其呈现”的整体作用作为评价对象，并在研究主张和分析中保持一致。

正式评价比较研究界面和基础界面在有依据的诊断达成率上的差异，并考察完成时间、材料核验、判断信心和使用负担。形成性研究确定设计机制后提出方向性假设，预实验完成后预注册具体假设、主要指标和评分规则。

\noindent\textbf{任务分配。}

正式实验的招募条件与形成性研究一致，并记录参与者经验。预实验拟招募4—6人，用于检查任务时长、理解难度、评分可操作性和测量流程。预实验参与者不进入正式实验；正式实验参与者不提前接触所用案例。

正式案例池初步按12个案例规划，每人完成6个，每种界面各3个。最终案例数量和单人任务量依据预实验中的任务时长和难度确定。通过平衡任务分配，使每个案例在两种界面下均获得评价，同一参与者只接触每个案例一次。界面顺序与案例分配交叉平衡，避免将某类故障固定给某一条件。全部任务时长和休息安排依据预实验确定。

样本量依据主要指标、最小有实际意义的差异、参与者和案例的变异范围进行模拟论证。Lakens\cite{ref67}的样本量论证方法以及 Green 和 MacLeod\cite{ref68}的模拟方法为此提供依据。研究在正式招募前确定目标人数和分析计划，将预实验参数与保守参数范围共同纳入计算，避免仅依赖单一试测效应。

\noindent\textbf{评价指标。}

每个任务要求参与者提交：对失败原因的判断、支持判断的运行事件、尚未排除的解释、下一步检查及判断信心。正式计时任务不要求持续口述，过程资料以操作记录为主，并在任务后进行简短回顾，减少口述对完成时间的影响。

主要指标为有依据的诊断达成率。每个答案依据预先确定的标准判断是否达成：

\begin{enumerate}
\item 原因或无法确定的判断符合案例可见材料允许的结论。
\item 引用的运行事件准确且能够支持关键判断。
\item 提出的下一步检查与剩余疑问相匹配；材料已经充分时，可提出有针对性的确认或修复检查。
\end{enumerate}

三个条件分别评分，再按预注册规则形成达成判定，同时保留分项结果。承认无法判断需要指出具体缺口和可区分解释的检查，空泛保留意见不计为达成。

\begingroup
\small
\setstretch{1.15}
\renewcommand{\arraystretch}{1.1}
\setlength{\LTleft}{0pt}
\setlength{\LTright}{0pt}
\begin{longtable}{@{}>{\raggedright\arraybackslash}p{0.22\dimexpr\textwidth-4\tabcolsep\relax}>{\raggedright\arraybackslash}p{0.39\dimexpr\textwidth-4\tabcolsep\relax}>{\raggedright\arraybackslash}p{0.39\dimexpr\textwidth-4\tabcolsep\relax}@{}}
\caption{诊断评价指标}\label{tab:selection-report-3}\\
\toprule
\textbf{指标类别} & \textbf{具体测量} & \textbf{解释用途} \\
\midrule
\endfirsthead
\multicolumn{3}{c}{表\thetable\quad 诊断评价指标（续）}\\
\toprule
\textbf{指标类别} & \textbf{具体测量} & \textbf{解释用途} \\
\midrule
\endhead
\bottomrule
\endfoot
主要结果 & 有依据的诊断达成率 & 比较两种界面对诊断任务的支持 \\
判断组成 & 原因或信息不足判断、引用依据、下一步检查的分项得分 & 识别改善发生的环节 \\
时间与操作 & 完成时间、关键材料访问、比较与回查行为 & 描述获得结果所需的成本及过程 \\
信心 & 作答后0—100分信心，与答案达成情况对应 & 检查错误判断是否伴随过高信心 \\
使用体验 & 心智努力、操作困难及访谈反馈 & 解释使用负担和接受程度 \\
\end{longtable}
\endgroup

任务达到时间上限时结束并保存当前判断；完成时间按未完成任务的截尾情况处理。两名评分者依据匿名化答案独立评分，评分材料隐藏界面条件，先报告一致性，再按评分手册讨论分歧。评分手册在正式数据分析前冻结。

\noindent\textbf{数据分析。}

主要结果拟使用广义线性混合模型（generalized linear mixed model），以界面条件为主要固定效应，纳入参与者与案例的随机效应，并按预注册计划处理顺序和经验因素。模型结构、参数估计及拟合异常的处理方式在预实验后确定。结果报告效应估计及置信区间。

时间、分项得分和使用体验分别分析。过程资料用于检查关键材料是否被发现、原因解释是否得到核验，以及错误判断如何形成。经验和案例条件的差异主要用于描述效果范围，探索性分析与主要检验分别呈现。

研究将主要结果、分项得分和过程资料与设计的预期作用对应分析。结论范围依据实际出现的诊断表现、检查行为和使用成本确定。

\subsubsection{迁移检验}

对照评价后，拟另外构建3—4个未参与设计迭代的新案例，改变任务内容、工具组合或框架实现中的至少一项条件，保留待检验的共同诊断问题。拟招募6—8名具有相关经验的参与者开展任务走查，分析同一设计操作在新案例中的使用过程、所需前提和失效情境。该阶段用于修订原则的适用范围，不承担新的总体效果检验。

研究同时尝试将一项核心原则用另一种低保真表现方式表达，并结合任务走查检查使用者能否完成相同检查活动。该步骤帮助区分原则依赖的操作与某个界面布局的偶然特征。结果将与正式实验一同用于修订原则。

每项原则按照以下结构整理：

\begingroup
\small
\setstretch{1.15}
\renewcommand{\arraystretch}{1.1}
\setlength{\LTleft}{0pt}
\setlength{\LTright}{0pt}
\begin{longtable}{@{}>{\raggedright\arraybackslash}p{0.24\dimexpr\textwidth-2\tabcolsep\relax}>{\raggedright\arraybackslash}p{0.76\dimexpr\textwidth-2\tabcolsep\relax}@{}}
\caption{设计原则的表述结构}\label{tab:selection-report-4}\\
\toprule
\textbf{内容} & \textbf{需要回答的问题} \\
\midrule
\endfirsthead
\multicolumn{2}{c}{表\thetable\quad 设计原则的表述结构（续）}\\
\toprule
\textbf{内容} & \textbf{需要回答的问题} \\
\midrule
\endhead
\bottomrule
\endfoot
适用情境 & 哪类任务、信息条件或使用者困难需要这项原则？ \\
设计操作 & 设计者应怎样组织信息或提供检查操作？ \\
预期作用 & 该操作影响哪项诊断活动？ \\
支持证据 & 哪些观察、比较结果和新案例支持这一判断？ \\
实施条件与成本 & 需要哪些记录、背景知识和操作投入？ \\
适用边界 & 哪些情况下效果减弱、无效或产生新的困难？ \\
\end{longtable}
\endgroup

研究保留未得到支持的设计主张，并说明其结果。最终原则数量根据证据决定。

\subsubsection{质量控制}

本研究重点是建立实践问题、设计选择与评价结论之间的对应关系。形成性研究确定可重复观察的困难，原型表达针对这一困难的设计主张，实验检查主张是否获得支持，新案例研究界定适用条件。

另一个重点是诊断答案的合理性。研究将区分实际故障、可见证据和参与者判断，避免只按最终原因名称评价任务表现。原始材料、派生信息和自动建议保留来源，使评分和过程分析具有可追溯依据。



\begingroup
\small
\setstretch{1.15}
\renewcommand{\arraystretch}{1.1}
\setlength{\LTleft}{0pt}
\setlength{\LTright}{0pt}
\begin{longtable}{@{}>{\raggedright\arraybackslash}p{0.24\dimexpr\textwidth-2\tabcolsep\relax}>{\raggedright\arraybackslash}p{0.76\dimexpr\textwidth-2\tabcolsep\relax}@{}}
\caption{研究风险控制}\label{tab:selection-report-5}\\
\toprule
\textbf{难点} & \textbf{应对安排} \\
\midrule
\endfirsthead
\multicolumn{2}{c}{表\thetable\quad 研究风险控制（续）}\\
\toprule
\textbf{难点} & \textbf{应对安排} \\
\midrule
\endhead
\bottomrule
\endfoot
实际案例包含公司或个人信息 & 以授权脱敏事件理解情境，用开源系统构建可复现任务 \\
自建案例与真实困难脱节 & 依据访谈选择案例，由实践者核查背景、歧义和检查路径 \\
多种原因符合表面症状 & 分别记录实际原因与可见材料支持的解释，允许多条合理诊断路径 \\
原型同时改变过多因素 & 围绕一个核心问题设计条件，核对各条件的信息和操作差异 \\
参与者经验与案例难度造成差异 & 记录经验，训练基础操作，平衡任务，联合分析人员与案例差异 \\
专业参与者招募不足 & 提前开展公司与外部渠道招募，分阶段维护招募矩阵并调整计划 \\
模型和工具更新影响重复评价 & 固定运行材料、版本与配置，保存构造记录及可复现环境 \\
\end{longtable}
\endgroup

\subsection{可行性分析}

课题来自实际开发工作中的问题观察，已完成文献调研和方案设计。访谈、案例、原型和实验尚未开展，文中样本及研究安排均属计划。

研究者可通过公司渠道接触若干具有智能体开发经验的人员，并计划通过外部专业社群扩展招募。技术实现拟使用开源框架，研究原型主要处理保存的运行材料和必要的检查操作，开发范围与核心研究问题保持一致。

整体工作按阶段推进。先以访谈和轻量案例确认问题，再通过低保真方案和任务走查确定设计，正式实现集中于评价需要的功能。中期前完成形成性研究、原型和预实验，并争取取得正式评价的阶段性数据；中期后完成评价、迁移检验和论文写作。

\subsection{研究伦理}

涉及参与者的研究将在完成学校要求的伦理程序后开展。知情同意材料说明研究目的、参与内容、录音及屏幕记录范围、补偿方式、退出权利和数据保存安排。公司渠道招募遵循自愿原则，参与情况不用于工作评价。

研究仅收集完成分析所需的材料。真实案例须取得授权并进行脱敏，排除密钥、个人信息及不可公开的业务内容。联系信息与研究编号分开保存，原始录音、屏幕记录和文本资料限制访问。论文使用匿名编号和必要摘录，公开材料遵守参与者授权及开源许可。

自动生成的摘要或建议保留生成配置与来源标记。案例、评分手册、原型版本和分析材料统一归档，以支持研究复核。

\section{研究框架}

根据 Zimmerman 等\cite{zimmerman2007rtd}关于通过设计开展研究的论述，以及 Meyer 和 Dykes\cite{meyer2020rigor}对设计研究严谨性的讨论，本研究依次开展诊断实践研究、设计探索、对照评价、新案例检验和知识提炼。各阶段的产出作为后续阶段的输入。

\begingroup
\small
\setstretch{1.15}
\renewcommand{\arraystretch}{1.1}
\setlength{\LTleft}{0pt}
\setlength{\LTright}{0pt}
\begin{longtable}{@{}>{\raggedright\arraybackslash}p{0.17\dimexpr\textwidth-6\tabcolsep\relax}>{\raggedright\arraybackslash}p{0.28\dimexpr\textwidth-6\tabcolsep\relax}>{\raggedright\arraybackslash}p{0.27\dimexpr\textwidth-6\tabcolsep\relax}>{\raggedright\arraybackslash}p{0.28\dimexpr\textwidth-6\tabcolsep\relax}@{}}
\caption{研究实施阶段}\label{tab:selection-report-2}\\
\toprule
\textbf{阶段} & \textbf{核心工作} & \textbf{阶段产出} & \textbf{进入下一阶段的依据} \\
\midrule
\endfirsthead
\multicolumn{4}{c}{表\thetable\quad 研究实施阶段（续）}\\
\toprule
\textbf{阶段} & \textbf{核心工作} & \textbf{阶段产出} & \textbf{进入下一阶段的依据} \\
\midrule
\endhead
\bottomrule
\endfoot
诊断实践研究 & 关键事件访谈、任务观察、跨案例分析 & 困难与信息条件矩阵 & 核心困难有多个案例支持，能够通过设计进行比较 \\
设计探索 & 提出备选方案、任务式走查、迭代原型 & 明确设计主张与原型 & 方案可用，设计差异可描述，评价材料能够回答问题 \\
对照评价 & 冻结条件、开展正式任务、分析判断和过程 & 效果估计与行为解释 & 获得足以讨论核心主张的任务及过程资料 \\
新案例检验 & 更换任务或工具条件，观察作用变化 & 支持案例、失效案例及修订依据 & 能够说明原则适用与不适用的条件 \\
知识提炼 & 综合实证材料与设计记录 & 经验描述和设计原则 & 原则各项论断均有对应材料与证据范围 \\
\end{longtable}
\endgroup

核心困难的选择依据包括出现范围、诊断后果、材料可获得性、设计可介入性和实验可检验性。原型围绕最终选定的问题实现必要功能。

\section{计划进度}

\begingroup
\small
\setstretch{1.15}
\renewcommand{\arraystretch}{1.1}
\setlength{\LTleft}{0pt}
\setlength{\LTright}{0pt}
\begin{longtable}{@{}>{\raggedright\arraybackslash}p{0.22\dimexpr\textwidth-4\tabcolsep\relax}>{\raggedright\arraybackslash}p{0.39\dimexpr\textwidth-4\tabcolsep\relax}>{\raggedright\arraybackslash}p{0.39\dimexpr\textwidth-4\tabcolsep\relax}@{}}
\caption{研究计划}\label{tab:selection-report-6}\\
\toprule
\textbf{时间} & \textbf{研究工作} & \textbf{阶段成果} \\
\midrule
\endfirsthead
\multicolumn{3}{c}{表\thetable\quad 研究计划（续）}\\
\toprule
\textbf{时间} & \textbf{研究工作} & \textbf{阶段成果} \\
\midrule
\endhead
\bottomrule
\endfoot
2026.09—2026.10 & 完成开题，细化研究协议，办理伦理程序，构建初始案例 & 研究协议、招募材料、案例模板 \\
2026.11—2027.01 & 开展关键事件访谈与任务观察，同步整理资料 & 访谈及观察资料、初步编码 \\
2027.02—2027.03 & 跨案例分析，补充样本，确定核心设计问题 & 困难与信息条件矩阵、设计要求 \\
2027.03—2027.05 & 比较备选方案，完成两轮任务式走查 & 设计记录、核心主张、原型方案 \\
2027.05—2027.06 & 实现研究原型，核查案例与评分规则 & 原型、正式案例池、评分手册 \\
2027.06—2027.08 & 预实验、样本量论证与预注册，启动正式评价 & 预实验结果、正式研究计划、阶段数据 \\
2027.09 & 中期答辩 & 形成性研究、设计依据、原型及评价进展 \\
2027.09—2027.10 & 完成正式评价与新案例检验 & 完整研究数据、原则修订依据 \\
2027.11—2028.02 & 综合分析并撰写论文 & 经验描述、设计原则、论文初稿 \\
2028.03—2028.05 & 论文修改、送审与答辩准备 & 学位论文定稿、归档材料 \\
\end{longtable}
\endgroup
